\chapter{METHODOLOGY}
\label{chap:method}

\section{System Overview}
\label{sec:overview}
The study proceeds in stages, shown in Figure~\ref{fig:pipeline}. A data generator produces group word problems whose input alphabet, and therefore whose per-token requirement, is controlled precisely. A single DeltaProduct layer processes each sequence and a linear readout predicts, after every token, the group element obtained so far. The layer either applies a fixed number of Householder factors to every token, receives its allocation from a reference derived from the theoretical analysis, or learns its allocation through a halting gate trained with a compute-budget penalty.

\begin{figure}[ht]
\centering
\resizebox{0.88\linewidth}{!}{%
\begin{tikzpicture}[
  box/.style={draw, rounded corners=2pt, align=center, minimum height=1.6cm, text width=3.1cm, font=\small, fill=blue!5},
  th/.style={draw, rounded corners=2pt, align=center, minimum height=1.3cm, text width=3.1cm, font=\small, fill=orange!10},
  arr/.style={-{Stealth[length=2.4mm]}, thick},
  darr/.style={-{Stealth[length=2.4mm]}, thick, dashed},
  lab/.style={font=\footnotesize, fill=white, inner sep=1pt}]
\node[box] (data)  at (0,0)    {Group word-problem data\\[2pt] alphabet $\Sigma$, mixture rate $p$};
\node[box] (layer) at (4.4,0)  {Linear RNN layer\\[2pt] Householder factors + halting gate};
\node[box] (read)  at (8.8,0)  {Readout\\[2pt] group element after every token};
\node[box] (eval)  at (13.2,0) {Evaluation\\[2pt] accuracy and allocated order};
\draw[arr] (data) -- (layer);
\draw[arr] (layer) -- (read);
\draw[arr] (read) -- (eval);
\node[th] (theory) at (0,-3.0)   {Order-requirement analysis $D(g)$};
\node[th] (budget) at (4.4,-3.0) {Compute-budget penalty $\gamma\,\bar n$};
\node[th] (oracle) at (8.8,-3.0) {Reference allocation (oracle)};
\draw[arr] (budget) -- node[lab, left, pos=0.55] {trains gate} (layer);
\draw[darr] (oracle.north) -- ++(0,0.75) coordinate (o1) -| ([xshift=7mm]layer.south);
\node[lab, above] at ($(o1)+(-1.3,0)$) {replaces gate};
\draw[darr] (theory.south) -- ++(0,-0.6) -| node[lab, pos=0.25, below] {per-token requirement} (oracle.south);
\draw[darr] (theory.south) ++(-0.5,0) -- ++(0,-1.25) -| node[lab, pos=0.2, below] {demand labels for allocation analysis} (eval.south);
\end{tikzpicture}}
\caption[Overview of the study]{Overview of the study. Solid arrows are the model and its training signal; dashed arrows carry information from the theoretical analysis, which is used only for the reference allocation and for evaluation, never as a training target for the learned gate.}
\label{fig:pipeline}
\end{figure}

The study is guided by three hypotheses:
\begin{itemize}[itemsep=1pt,topsep=3pt]
	\item \textbf{H1:} On inputs whose tokens differ in the expressivity they require, an allocation that follows each token's requirement reaches the accuracy of the maximal fixed order at a lower required order per token.
	\item \textbf{H2:} The minimal number of Householder factors a token requires is determined by the group generated by the entire input alphabet, not by the token alone.
	\item \textbf{H3:} A monotone halting gate trained end to end with a compute-budget penalty can recover the per-token requirement without supervision on the allocation.
\end{itemize}

\section{Task Formulation: Group Word Problems}
\label{sec:task}
Five items occupy five positions (Figure~\ref{fig:task}). The model reads a sequence of instructions, one per token, each of which rearranges the items, and after every instruction it must output the arrangement reached so far. Formally, each token $x_t$ is a permutation in a finite group $G$, and the target at position $t$ is the running product
\begin{equation}
y_t = x_t\, x_{t-1} \cdots x_1 .
\label{eq:target}
\end{equation}
Because $y_t$ depends on every earlier token, the task cannot be solved by inspecting any single token, and because the answer is known exactly, accuracy is unambiguous.

\begin{figure}[ht]
\centering
\resizebox{0.97\linewidth}{!}{%
\begin{tikzpicture}[
  cell/.style={draw, minimum size=0.78cm, font=\normalsize},
  mv/.style={cell, fill=orange!18},
  arr/.style={-{Stealth[length=2.6mm]}, very thick},
  cp/.style={font=\small, align=center}]
\foreach \v [count=\i] in {1,2,3,4,5} { \node[cell] at (\i*0.78, 0) {\v}; }
\node[cp] at (2.34, 0.85) {initial state};
\node[cp] at (2.34, -0.85) {$y_0 = e$};
\draw[arr] (4.55,0) -- node[above=2pt, cp] {$x_1$: swap 2 and 4} node[below=2pt, cp] {$\ell(x_1) = 1$} (7.65,0);
\foreach \v [count=\i] in {1,4,3,2,5} {
  \ifnum\i=2 \node[mv] at (7.52+\i*0.78, 0) {\v};
  \else\ifnum\i=4 \node[mv] at (7.52+\i*0.78, 0) {\v};
  \else \node[cell] at (7.52+\i*0.78, 0) {\v}; \fi\fi }
\node[cp] at (9.86, 0.85) {after token 1};
\node[cp] at (9.86, -0.85) {$y_1 = x_1$};
\draw[arr] (12.07,0) -- node[above=2pt, cp] {$x_2$: 5-cycle\\ $(1\,2\,3\,4\,5)$} node[below=2pt, cp] {$\ell(x_2) = 4$} (15.17,0);
\foreach \v [count=\i] in {5,1,4,3,2} { \node[mv] at (15.04+\i*0.78, 0) {\v}; }
\node[cp] at (17.38, 0.85) {after token 2};
\node[cp] at (17.38, -0.85) {$y_2 = x_2 x_1$};
\end{tikzpicture}}
\caption[The group word problem on five items]{The group word problem on five items. Each token is a permutation; after every token the model must output the arrangement reached so far. Shaded cells are those moved by the most recent token. The transposition length $\ell$ counts the fewest swaps that build each token.}
\label{fig:task}
\end{figure}

The \emph{symmetric group} $S_n$ contains all $n!$ rearrangements of $n$ items. A \emph{transposition} swaps two items and a \emph{5-cycle} moves all five around one loop. The fewest transpositions that compose to $g$ is its \emph{transposition length} $\ell(g) = n - c(g)$, with $c(g)$ the number of cycles, so a transposition has $\ell=1$ and a 5-cycle $\ell=4$. The \emph{even} permutations (even $\ell$) form the \emph{alternating group} $A_n$, half of $S_n$; a 5-cycle is even and a transposition odd.

An \emph{alphabet} $\Sigma$ is the set of permutations the generator may emit, and the group the model must track is its \emph{closure} $\gen{\Sigma}$, the set of all products of alphabet elements. An alphabet of 5-cycles contains only even permutations, so its closure is $A_5$; adding a single transposition makes odd arrangements reachable and the closure becomes $S_5$. Table~\ref{tab:alphabets} lists the alphabets used. In the \emph{mixture} alphabets each token is a 5-cycle with probability $p$ and a transposition otherwise; the \emph{closure} alphabets are sampled uniformly and are paired so that they share the same 5-cycles but differ in their closure. Sequences have 128 tokens; the output vocabulary has 127 classes (the 120 elements of $S_5$ and 7 special tokens).

\begin{table}[ht]
	\caption{Alphabets used in the experiments}
	\centering
	\footnotesize
	\setlength{\tabcolsep}{4pt}
	\begin{tabular}{p{2.5cm} p{4.4cm} c c c}
		\toprule
		\textbf{Alphabet} & \textbf{Contents} & \textbf{Symbols} & $\gen{\Sigma}$ & \textbf{Share at $D_{\max}$} \\
		\midrule
		Mixture, rate $p$ & 10 transpositions, 24 five-cycles; five-cycle with probability $p$ & 34 & $S_5$ ($p<1$) & $p$ \\
		\cell{A5\_c5}     & 24 five-cycles                     & 24 & $A_5$ & 1.00 \\
		\cell{A5\_c5\_dt} & five-cycles, double transpositions & 39 & $A_5$ & 1.00 \\
		\cell{A5\_c5\_3c} & five-cycles, 3-cycles              & 44 & $A_5$ & 1.00 \\
		\cell{S5\_c5\_t1} & five-cycles, \emph{one} transposition & 25 & $S_5$ & 0.96 \\
		\cell{S5\_c5\_t}  & five-cycles, all 10 transpositions & 34 & $S_5$ & 0.71 \\
		\bottomrule
	\end{tabular}
	\label{tab:alphabets}
\end{table}

\section{Base Architecture}
\label{sec:base}
The model is a single DeltaProduct layer \cite{deltaproduct} between a token embedding and a linear readout. For each token the layer computes keys $k_{t,i}$, values $v_{t,i}$ and strengths
\begin{equation}
\beta_{t,i} = 2\,\sigma\!\left(w_i^\top h_t\right) \in (0,2), \qquad i = 1,\dots,\nh,
\label{eq:beta}
\end{equation}
where $h_t$ is the token representation, and applies $\nh$ consecutive delta-rule steps (Equation~\ref{eq:deltarule}), giving the transition of Equation~\ref{eq:deltaproduct}; the range $(0,2)$ lets each factor act as anything from the identity to a reflection \cite{negeig}. The implementation uses the gated DeltaProduct layer of \texttt{flash-linear-attention} \cite{fla}, and the readout is trained with cross-entropy over all positions.

\section{Minimal Order Requirement of a Token}
\label{sec:demand}
\begin{proposition}[Rank bound]
\label{prop:rank}
For any strengths $\beta_i$ and unit vectors $k_i$, $A = \prod_{i=1}^{K}(I - \beta_i k_i k_i^\top)$ satisfies $\rank(A - I) \le K$.
\end{proposition}
\begin{proof}
$H_1\cdots H_m - I = (H_1\cdots H_{m-1} - I)H_m + (H_m - I)$, and $H_m - I$ has rank at most one, so each factor raises $\rank(\cdot - I)$ by at most one.
\end{proof}
For a permutation matrix, $\rank(P_g - I) = n - c(g) = \ell(g)$, so a layer that realises token $g$ as its permutation matrix needs at least $\ell(g)$ factors, whatever the strengths.

\begin{proposition}[Parity bound]
\label{prop:parity}
If every strength is fixed at $\beta_i = 2$, then $\det A = (-1)^K$ and $P_g$ is reachable only if $K \ge \ell(g)$ and $K \equiv \ell(g) \pmod 2$. With strengths free in $[0,2]$, every $K \ge \ell(g)$ suffices, since surplus factors can be set to $\beta = 0$, the identity.
\end{proposition}
Proposition~\ref{prop:parity} fixes a design decision: to use fewer than its maximum number of factors, the layer must switch a factor off exactly, which requires $\beta=0$. The halting gate therefore \emph{scales} factor strengths rather than masking factors of fixed strength.

\textbf{The role of the representation.} A \emph{representation} $\rho$ of a group maps each element to an invertible matrix with $\rho(gh) = \rho(g)\rho(h)$; it is \emph{faithful} if distinct elements map to distinct matrices. Permutation matrices are one representation, but a group generally has others of different dimensions. A layer need only keep group elements distinguishable, so it may use any faithful representation, and Proposition~\ref{prop:rank} then applies to $\rho(g)$; for orthogonal representations the bound is attained, since by the Cartan--Dieudonn\'e theorem an orthogonal map is a product of $\rank(\rho(g) - I)$ reflections \cite{geometricalgebra}.

\begin{definition}[Order requirement]
\label{def:demand}
For an alphabet $\Sigma$ and a token $g \in \Sigma$,
\begin{equation}
D(g) \;=\; \min_{\rho\ \text{faithful rep.\ of}\ \gen{\Sigma}} \ \rank\!\left(\rho(g) - I\right).
\label{eq:demand}
\end{equation}
\end{definition}

The minimum ranges over representations of the closure, not of the token. For alphabets generating $S_5$, the four-dimensional standard representation gives $D(g) = \ell(g)$ and no faithful representation does better (Section~\ref{sec:res-verify}). $A_5$, however, is isomorphic to the rotation group of the icosahedron, a faithful three-dimensional representation that $S_5$ lacks. Every non-identity rotation fixes exactly its axis, so $\rank(R - I) = 2$ for every non-identity element of $A_5$ (Table~\ref{tab:a5classes}).

\begin{table}[ht]
	\caption{Order requirement of the element classes of $A_5$}
	\centering
	\begin{tabular}{l c c c}
		\toprule
		\textbf{Element class} & \textbf{Count} & \textbf{$\ell$ (perm.\ rep.)} & \textbf{Rank in 3-D rep.} \\
		\midrule
		Identity             & 1  & 0 & 0 \\
		Double transposition & 15 & 2 & 2 \\
		3-cycle              & 20 & 2 & 2 \\
		5-cycle              & 24 & 4 & \textbf{2} \\
		\bottomrule
	\end{tabular}
	\label{tab:a5classes}
\end{table}

Hence the requirement of a token is a property of the whole alphabet: adding one transposition to an alphabet of 5-cycles changes the closure from $A_5$ to $S_5$ and raises the requirement of every 5-cycle from two to four, which no model of per-token cost can produce (tested in Section~\ref{sec:res-closure}). Per-token allocation can also only help where $D$ varies; for closures $A_4$ and $A_5$ it does not (Table~\ref{tab:groups}, Appendix~I).

\section{The Adaptive-Order Layer}
\label{sec:layer}
The layer keeps a maximum order $K$ and adds a halting gate, a two-layer perceptron of hidden width 64 applied to the token representation, $z_t = W_2\,\mathrm{GELU}(W_1 h_t + b_1) + b_2 \in \mathbb{R}^K$. Its outputs become gates through a cumulative product:
\begin{equation}
\lambda_{t,i} = \sigma(z_{t,i}), \qquad
g_{t,i} = \prod_{j=1}^{i} \lambda_{t,j}, \qquad
n_t = \sum_{i=1}^{K} g_{t,i}, \qquad
\tilde\beta_{t,i} = g_{t,i}\, \beta_{t,i}.
\label{eq:gates}
\end{equation}
The gates are monotone, $g_{t,1} \ge \cdots \ge g_{t,K}$, so once the gate closes it stays closed, and $n_t$ is the \emph{effective order} of token $t$. Following Proposition~\ref{prop:parity}, the gate scales each strength; a factor with $g_{t,i} = 0$ erases and writes nothing and is exactly the identity. The layer is shown in Figure~\ref{fig:layer}; its forward pass is given as Algorithm~\ref{alg:forward} in Appendix~I. During training the gates are real-valued; at evaluation they are rounded, $\hat g_{t,i} = \mathbf{1}[g_{t,i} > 0.5]$, so that the reported order is an integer number of factors actually applied.

\begin{figure}[ht]
\centering
\resizebox{0.56\linewidth}{!}{%
\begin{tikzpicture}[
  b/.style={draw, rounded corners=2pt, align=center, text width=4.1cm, minimum height=1.05cm, font=\small, fill=blue!5},
  g/.style={b, fill=green!8},
  p/.style={b, fill=orange!10, text width=3.2cm},
  arr/.style={-{Stealth[length=2.3mm]}, thick},
  darr/.style={-{Stealth[length=2.3mm]}, thick, dashed}]
\node[b] (h) at (0,0) {Token representation $h_t$};
\node[g] (mlp) at (-3.3,-1.8) {Halting gate MLP\\ $128 \to 64 \to K$, GELU};
\node[g] (sig) at (-3.3,-3.4) {$\lambda_{t,i} = \sigma(z_{t,i})$ in float32};
\node[g] (cum) at (-3.3,-5.0) {$g_{t,i} = \prod_{j\le i}\lambda_{t,j}$,\quad $n_t = \sum_i g_{t,i}$};
\node[b] (proj) at (3.3,-1.8) {Key, value and strength projections};
\node[b] (bet) at (3.3,-3.4) {$\beta_{t,i} = 2\sigma(w_i^\top h_t)$,\\ $k_{t,i}$, $v_{t,i}$, \ $i = 1..K$};
\node[b] (sc) at (0,-6.7) {Gated strengths\\ $\tilde\beta_{t,i} = g_{t,i}\,\beta_{t,i}$};
\node[b] (prod) at (0,-8.4) {$K$ delta-rule steps:\\ $A_t = \prod_{i}(I - \tilde\beta_{t,i}k_{t,i}k_{t,i}^\top)$};
\node[b] (st) at (0,-10.0) {State update $S_{t-1} \to S_t$, readout};
\node[p] (bud) at (-7.0,-6.7) {Budget penalty\\ $\gamma\cdot\frac{1}{BT}\sum_t n_t$};
\draw[arr] (h) -- (mlp); \draw[arr] (h) -- (proj);
\draw[arr] (mlp) -- (sig); \draw[arr] (sig) -- (cum);
\draw[arr] (proj) -- (bet);
\draw[arr] (cum) -- (sc); \draw[arr] (bet) -- (sc);
\draw[arr] (sc) -- (prod); \draw[arr] (prod) -- (st);
\draw[darr] (cum.west) -| (bud.north);
\end{tikzpicture}}
\caption[The adaptive-order layer]{The adaptive-order layer. The halting gate (green) emits monotone gates that scale the strength of each Householder factor; a factor whose gate is zero is the identity. The effective order $n_t$ enters the loss through the budget penalty (orange) during training.}
\label{fig:layer}
\end{figure}



\textbf{Training objective.} Nothing in the task loss rewards using fewer factors, so a penalty on the effective order is added,
\begin{equation}
\mathcal{L} = \mathcal{L}_{\text{task}} + \gamma(s)\cdot \frac{1}{BT}\sum_{b,t} n_{b,t}, \qquad
\gamma(s) = \gamma \cdot \min\!\left(1,\ \max\!\left(0,\ \frac{s - s_w}{s_r}\right)\right),
\label{eq:loss}
\end{equation}
where $s$ is the training step. The penalty is off for the first $s_w = 4000$ steps, so that the model first learns the task, and then rises linearly over $s_r = 4000$ steps to its final weight $\gamma$ (Algorithm~\ref{alg:train} in Appendix~I).

\textbf{Initialisation and precision.} The gate's output layer starts with zero weights and bias $4$, so the initial effective order is $\sum_{i=1}^{4}\sigma(4)^{i} = 3.8233$ for every token: the gate starts open and must learn to close. The recurrence kernels run in bfloat16, in which a sigmoid rounds to exactly $1.0$, with zero gradient, for every logit above $6.2364$; the gate is therefore evaluated in float32, which moves this threshold to $17.33$ (Table~\ref{tab:precision}, Appendix~I).

\section{Reference Allocation}
\label{sec:oracle}
To measure what exact per-token allocation achieves independently of whether it can be learned, the gate is replaced by a \emph{reference allocation} (the oracle), $g^{\text{oracle}}_{t,i} = \mathbf{1}[\, i \le \ell(x_t) \,]$. For alphabets generating $S_5$ this equals $D(x_t)$. The oracle uses the same layer and training procedure as the learned model; only the source of the gates differs. It reads the target structure and serves only as a reference.
